\begin{abstract}

Short-text classification in Portuguese faces three obstacles at once. The vocabulary is abbreviated and noisy, as in product descriptions from electronic invoices. The categories number in the hundreds, and are imbalanced. And the method depends on large volumes of labels, whose acquisition is the dominant cost of these systems. This thesis proposes and evaluates FALCO (Active Learning Framework with LLMs for Short Text), which reduces that cost in three stages. The first is informed composition of the initial training set. The second is uncertainty-based active selection of instances. The third is partial replacement of the human annotator by oracles based on large language models (LLMs), organized in phases of increasing cost. The investigation uses an audited public dataset of 250,365 product descriptions across 621 categories, of which 250,221 in the corrected version used in the experiments. The protocol is pre-registered, with immutable partitions, paired analysis (McNemar, Wilcoxon), and a traceable artifact for every reported number. Five results support the thesis. (i) The composition of the initial set matters more the smaller the budget: random draws alone produce a $6.4$ percentage-point accuracy range at $|L_0|=100$. The proposed deterministic heuristic, DRI-SL, combines semantic density and lexical diversity while using no labels at all. It outperforms the best individual found by a supervised genetic algorithm at every size from 100 to 5,000. The comparison is conservative, since the circularity correction introduced in this thesis reduced the evolutionary envelope by up to $6.3$ percentage points. (ii) Zero-shot LLM oracles reach $77$--$83\%$ accuracy in the closed space of 621 categories, at costs of US\$~0.035 to US\$~0.92 per thousand labels. Their Macro F1 exceeds that of a light supervised classifier trained on the full base of 250 thousand labeled descriptions. Nearly half of their errors concentrate on catalog conventions: sibling categories and umbrella classes. (iii) The measurement instrument is part of the result. Without output restriction to the class space, $6.8\%$ of responses are counted as spurious errors. Boundary rules in the prompt improve the weak model on the production distribution ($+4.6$ p.p., $p=0.012$), at the cost of severe degradation on rare classes ($p<0.001$). The serving provider is an instrument too: the same free model diverges significantly from itself across two services ($p<0.001$). Properly served, however, it ties with paid oracles on the production sample. (iv) Uncertainty strategies outperform random selection at the maximum significance achievable with 8 seeds ($p=0.0078$), and recover $78\%$ of the supervised ceiling's Macro F1 with $15\%$ of that experiment's pool. The advantage survives intact under uniform oracle noise up to $\varepsilon=0.4$, a range that covers the error observed in the real LLMs. No oracle reached the pre-registered $85\%$ accuracy criterion. The methodological gate therefore defined the final configuration: an economical oracle in the initial phase and the most accurate oracle in the advanced phase. It also made learning with noisy labels the central scenario. (v) The framework was executed end to end with a free LLM oracle, at zero monetary cost. Validation used the strong classifier, BERTimbau, with three seeds and evaluation on the reserved population of 177,490 instances. It answered the central hypothesis: obtaining at least $95\%$ of the accuracy of full supervision on the reference pool, the pre-registered metric, with at most 34,724 labels ($15\%$ of the deduplicated population). The verdict is twofold. \textbf{The criterion is attainable within the ceiling, and is crossed at 20,000 labels ($8.6\%$ of the population) in all three seeds, in a \textit{post hoc} gold-label sweep. The executed configuration, however, stopped by stagnation short of the floor, at 11,936 labels ($5.2\%$)}. The paired decomposition does not attribute the loss to oracle noise, which, being structured, even benefits rare classes. Nor does it attribute the loss to selection: it comes from the stopping criterion inherited from the light classifier. In Macro F1, the robustness analysis introduced in this thesis, the criterion is met at 30,000 labels ($13.0\%$ of the population), also in all three seeds and within the ceiling. And, on the average of the three seeds, 35,000 actively selected labels surpass full-pool supervision in both metrics, with heterogeneity declared: two seeds keep a significant advantage, one inverts in Macro F1 and ties in accuracy. Population-scale experiments further show that self-evaluation on actively collected data is biased in a metric-dependent direction: self-evaluation accuracy underestimates real performance, and self-evaluation Macro F1 overestimates it. This grounds the reserved evaluation set and the release criterion the framework adopts. The first-order limitations are declared. The results come from a single dataset, from the author's own study domain. The 35,000-label arm's edge over full supervision holds on the seed average, with one divergent seed in Macro F1. And the cost analysis rests on ratios between oracles, such as the 26-fold difference under a statistical accuracy tie, rather than on absolute prices, which date quickly. The thesis also delivers the LCE (Learning Curve Efficiency) metric, for learning-curve comparison, the open-source \texttt{activelearning} library, and the FlowBuilder tool. With sanitized data ingestion, parameterized execution, and learning curves, they make the protocol reusable in other domains.

\end{abstract}
